% 1. Preamble
\documentclass[12pt, a4paper]{article}

% Import Packages
\usepackage[utf8]{inputenc}        % Encoding
\usepackage{amsmath, amssymb}     % Advanced math formulas
\usepackage{mathtools}           % Math enhancements
\usepackage{graphicx}            % Image insertion
\usepackage{hyperref}            % Interactive hyperlinks & PDF bookmarks
\usepackage{color}               % Color text and backgrounds
\usepackage{geometry}            % Page geometry
\usepackage{comment}             % Multi-line comments
\usepackage{titling}             % Custom title formatting

\geometry{a4paper, margin=1in}

% Custom Command Definitions
\newcommand{\Sball}{\mathcal{S}}
\providecommand{\norm}[1]{\left\lVert#1\right\rVert}

% Hyperref Setup for Clean PDF Links
\hypersetup{
    colorlinks=true,
    linkcolor=blue,
    urlcolor=blue,
    pdftitle={Pattern Arithmetic: Higher Truths from Assembly},
    pdfauthor={Sandeep Lakshminarayan Chiluveru}
}

% Metadata & Custom Title Banner
\pretitle{\begin{center}\Huge\fontfamily{tt}\selectfont ================================================================================\\[0.5em]}
\posttitle{\\[0.5em] ================================================================================\end{center}}

\title{Pattern Arithmetic: Higher Truths from Assembly}
\author{Sandeep Lakshminarayan Chiluveru}
\date{\today}

\begin{document}

\maketitle

% ABSTRACT
\begin{abstract}
This is a second companion paper to \emph{Pattern Arithmetic: Typed
Algebra of Bounded Units on a Phase-Colored Sphere}, alongside
\emph{Pattern Arithmetic: Inference from Bowls and Paths}. The first
defines assembly ($\bowtie$) and an enclosure in which units may
descend an energy and rest as one camp or several; the second reads
a bowl of evidence into a verdict, but never asks what a settled
camp itself is entitled to be. This paper collects what is currently
known about that question, and it is deliberately small. Two points
are already closed, as direct consequences of choices the companion
paper already made rather than of anything built here: the
actor/patient question, settled at the level of a label rather than
a difference in motion, and the absence of astrophysical mass
segregation, a forced consequence of excluding volume from motion.
One point is genuinely open --- fine-grained multi-vector matching,
where an atom's own bundle, not its coarse derived vector, decides
how it bonds. What this paper does not yet do is the harder thing
its own title promises: say what a frozen camp, read and re-read at
successive checkpoints, is entitled to be called, and whether that
naming is itself a higher unit of truth or only a description of
one. That is stated here as this paper's own direction, not
answered as its content.
\end{abstract}

\section{Introduction}
\label{sec:intro}

A camp settles. Something froze at a resolution $\eta$, with a body
$\vec P_{c}$ and a membership --- and the companion paper is
explicit that the body is a description, measured after the fact,
never a destination the motion sought. Read again at a finer $\eta$,
or later in time, the same units may freeze into a different shape,
or one pile where there were two: a camp is a checkpoint, not a
fixed thing, and two checkpoints close enough to each other are
read as the same freeze only by a further, named allowance
(companion paper, Enclosure). None of this says what a checkpoint
\emph{is}, in the sense a bowl of evidence becomes a verdict in the
other companion paper. That is the question this paper exists to
carry, separately from inference, so that it can grow on its own
rather than as an appendix to a paper about something else.

What follows is not yet that answer. It is the ground already
cleared: two questions the companion paper's own axioms settle as a
side effect, stated here plainly since neither was recorded at the
time, and one question about how an atom's inner structure governs
assembly that remains genuinely open. The larger question --- what a
settled camp, read against known patterns or against its own
history of checkpoints, is entitled to be named, and whether that
naming closes into a symbol a higher unit could itself be built
from --- is named in the abstract and left there. Early exploration
toward it (reading a settled camp against known references, and
characterizing when such a reading is trustworthy rather than
lucky) exists only as worked, unpublished calculation; none of it is
claimed here as a result.

\section{What Is Inherited}
\label{sec:inherited}

We use, without re-deriving, the following from the companion
paper. Equation and axiom numbers refer to that paper.

\begin{itemize}
\item The atom $x=(r,\hat n,\chi,\tau)\in\mathcal{S}$, with $r\le 1$
      a normalized strength (Axiom 1).
\item Assembly, $\bowtie$: two units related while both remain
      intact, $A\bowtie B=(A,B,\kappa_{AB},U_{AB})$, with
      $U_{AB}<0$ for nonzero strengths and never repelling
      (Axiom~6).
\item The enclosure: external seats $\vec P$ descend the energy
      $\sum U_{ij}$ without mass or momentum until a stated force
      threshold $\eta$ is crossed; a configuration below it is a
      \emph{checkpoint}, not an exact equilibrium, since the pull
      between any two units of nonzero strength never reaches zero
      at finite distance (Enclosure, Rest and Camps).
\item Informational vectors (Axiom~7): an atom may carry $N$ capped
      arrows and one uncapped derived sum $V_{\mathrm{derived}}$.
      Coarse assembly runs on the derived pair; the bundle's own
      arrows are not consulted by $U_{AB}$ unless named explicitly.
\item \emph{Levels are closed}: a camp carries no representation of
      what it constitutes, and the enclosure supplies no comparison
      by which a checkpoint could name itself.
\end{itemize}

\section{What Assembly Already Settles}
\label{sec:already_settled}

Two settled points are worth recording plainly, since both follow
from Axiom~V4's exclusion of $\mathcal{V}_{\text{atom}}$ from motion
(companion paper, Enclosure \S~Rest) rather than from anything built
here.

\paragraph{The actor/patient question.} Closed at the level of a
label: polarity ($\lvert r_A-r_B\rvert$) distinguishes a mutual bond
from a one-sided one, and no further behavioral mechanism is
missing --- one was tested directly and found to require exactly
the volume-weighted motion that exclusion rules out on stability
grounds, so the honest final scope of this question is a
description of the bond, not a difference in how it moves.

\paragraph{Mass segregation.} The astrophysical phenomenon ---
stronger members sinking toward a group's centre --- was named
early in the companion paper's own development as an expected
consequence, while its dynamics was still built around a mass-like
quantity. Once motion settled as excluding $\mathcal{V}_{\text{atom}}$
entirely, that mechanism no longer applies here, confirmed directly:
equal force with nothing dividing it gives identical movement
regardless of strength or volume. The astrophysical phenomenon keeps
its real name; nothing in this framework currently reproduces it,
since the quantity that would drive it was deliberately kept out of
motion.

\section{Open: Fine-Grained Multi-Vector Matching}
\label{sec:open_matching}

Coarse assembly runs on the derived vector (companion paper,
Axiom~7), so $U_{AB}$ never needed to generalize; simultaneous
multi-partner bonding via informational vectors directly remains
optional, with an unconstrained double sum and an exclusive
(Hungarian) matching identified as the two candidates. They agree
whenever alignment quality is roughly uniform across candidates, but
rank-order pairing can be strictly worse than optimal when strength
and alignment pull in different directions --- verified with a
constructed adversarial case ($-1.17$ for rank order versus $-1.68$
true optimum).

\section{Scope and the Question This Paper Is For}
\label{sec:scope}

This paper is small on purpose. It records what the companion
paper's own choices already settle about assembly, once, so that
neither has to be rediscovered; and it names one open question about
how an atom's own bundle should govern multi-partner bonding, where
the coarse derived-vector treatment is known to be incomplete.

What it does not do is name a checkpoint. A frozen camp has a body
and a membership and, by the companion paper's own accounting,
nothing else --- no shape, no representation of what produced it,
no claim on what it should be called. Whether a checkpoint can be
read against known patterns and trusted, whether a sequence of
checkpoints under an allowance converges on something worth naming,
and whether such a naming is a genuine higher unit of truth or only
a compact description of a lower one --- these are the paper's own
stated direction, not its content. They are left here, deliberately,
for the work this paper exists to carry.

\section*{Development Notes: Ideas in Progress}
\addcontentsline{toc}{section}{Development Notes: Ideas in Progress}

What follows is a working record, not a finished section of this
paper. Each item is tagged by status --- \textbf{[Settled]},
\textbf{[Verified]}, \textbf{[Proposed]}, or \textbf{[Open]} --- so a
later pass can tell at a glance what is actually established, what
is argued but unbuilt, and what has been tried and genuinely not
yet solved, without re-deriving that distinction from scratch.
Nothing here is claimed as this paper's content; it is the material
its later sections would be built from, once built.

\subsection*{Structure and organization}

\textbf{[Proposed.]} Assembly as a layer beneath Inference, not a
peer to it: Inference reads whatever atoms already exist; Assembly
asks what happens when atoms relate to each other before anything
is read. A camp is not evidence, by the companion paper's own words
--- its body is ``measured afterwards\ldots a description rather
than a destination'' --- and this paper's own abstract already names
camp-naming as unanswered here, not in Inference. The layering is
implicit in language already sitting in both documents; stating it
as the organizing structure is what is new. Not yet written into
the Introduction.

\textbf{[Proposed.]} An eight-part restructuring built around that
layering: a reframed Introduction stating the one-directional
dependency; \emph{What a Camp Is}, an axiom-like core of invariant
camp properties; \emph{Two Modes --- Free and Directed Assembly};
\emph{Representing a Camp}; \emph{Structural Diagnostics}; the
existing inherited and open-question sections unchanged; a closing
statement rewritten around the layering. Two of eight sections need
no work; the other six, including the entire axiom-like core, are
still only conversation.

\textbf{[Proposed.]} Assembly's own closure axioms should carry
their own numbering, not continue the companion paper's --- justified
by that paper's own Axiom~0 discipline: a snowflake's native law is
recorded as refused a map onto the sphere rather than forced into
$\mathrm{SO}(3)$, because a system of axioms built for one kind of
question should not silently claim authority over a different kind
by mere continuation of its own numbering. Closure is that kind of
different question relative to the companion paper's atom-level
axioms.

\textbf{[Settled.]} The test for what actually earns axiom status
here --- forced (no freedom, like $\pi$), genuinely invented but
low-stakes (a normalization, like $r\le1$), or a real structural
choice (a different choice produces a different theory) --- is
already written into this paper, below, as the parked aside for the
companion paper's Thesis. Not repeated here.

\textbf{[Proposed.]} A first candidate that passes that test: what a
freeze keeps and what it discards. The companion paper says a
camp's body is a description, not a destination, and that a camp is
a checkpoint, not a destination, but never formalizes exactly what
survives a freeze versus what does not. Not derivable from the
companion paper's own axioms; not yet drafted as a formal statement.

\subsection*{Representing a camp}

\textbf{[Verified.]} Resultant versus dominant-vector representation,
tested on identical seats and bundles: $16.0\%$ of contested pairs
disagree ($120/749$ across 15 trials) on who ends up grouped with
whom, depending solely on which representation is used.

\textbf{[Verified.]} With strength made genuinely irrespective ---
removed not just from the pull's magnitude but from which vectors
get compared at all --- the two candidates already named in this
paper's own open question behave very differently. The unconstrained
double sum collapses to a single undifferentiated group in every
trial tested, confirmed not a settling-parameter artifact by a
same-population sanity check against the ordinary single-vector case.
Hungarian best-match survives with real, differentiated structure
(3--5 groups per trial, genuine size variation). Between this
paper's own two candidates, one of them is now shown not to work for
this case.

\textbf{[Proposed, untested.]} Centrality, treating $U_{AB}$ as a
weighted graph and reading a camp's most central member, as a fourth
candidate. Not yet tested at all.

\subsection*{Structural diagnostics}

\textbf{[Verified.]} Graph community detection is not a faster
substitute for physical settling. Tested properly, with distance
included in the edge weights after an initial pass without it was
caught as an unfair comparison: $69.7\%$ of contested pairs disagree,
and community detection consistently finds fewer, larger groups than
settling does. Diagnosis: settling captures the dynamic, cascading
resolution of competing pulls over time; a one-shot graph read of
the starting configuration cannot see that resolution happening.

\textbf{[Proposed, untested.]} Hierarchical clustering as the formal
version of a sentence the companion paper already states in words:
``a finer $\eta$ runs longer and may see groups merge that a coarser
one froze apart.'' That is a dendrogram, described in prose, never
built.

\subsection*{The query/selection border case}

\textbf{[Verified.]} A query matching two named vectors while
excluding a third produces a genuinely non-transitive relation ---
global transitivity checked directly at exactly zero, with individual
components confirmed to be chains rather than cliques (one
three-member component held only 2 of its 3 possible internal edges).
This is not a camp in this paper's sense; it is a sparse,
bridge-like network, closer to what sociology calls weak ties or
structural holes.

\textbf{[Flagged, unresolved.]} This looks structurally closer to
Inference's own open selection rule --- who enters, under what
condition --- than to Assembly's job of settling a membership once
given. Which paper should own it is not decided.

\subsection*{Systems as atoms}

\textbf{[Settled.]} Vocabulary: \emph{camp seat} for a camp's spatial
displacement (extending the companion paper's own existing ``seat''
register) versus \emph{camp signature} for its effective
informational reading (checked clean of collisions across all three
documents). These were previously, silently, one conflated vector.

\textbf{[Verified.]} Never-repels survives the jump to the system
level: checked directly on system-level pairs, minimum pairwise
coefficient positive ($\approx4.4\times10^{-7}$), nothing negative
anywhere.

\textbf{[Verified.]} The clip-versus-bijection choice matters again
at the system level, not only at the evidence level: the identical
16 camps produce 13 systems under the Hill bijection and 11 under
the naive clip, with different groupings, not merely different
counts.

\textbf{[Verified --- the hypothesis's central test.]} Treating
settled camps as new, system-level atoms produces genuinely different
behavior from simply widening the grouping radius on the original
atoms. Two-level construction: largest system holds $32\%$ of all
atoms ($16/50$). Naive coarsening of the same population at a
comparable scale: $94\%$ ($47/50$), a near-total collapse. Camps
compressing into a signature before the next level's dynamics run is
not equivalent to loosening the same single-level process.

\subsection*{Stability at the system level --- the open frontier}

\textbf{[Open, actively unsolved.]} Pure radial attraction with
nothing to counteract it trends toward full collapse --- matching the
companion paper's own stated Enclosure behavior, ``left longer, all
would merge into one,'' and matching why naive coarsening above jumps
almost directly to one giant group.

Checked against real astrophysics rather than assumed: globular
clusters do sometimes merge, but genuine mergers cluster around
specific circumstances --- early galactic epochs, dense environments,
close tidal encounters --- not a generic, ongoing process. Most
clusters, once in stable orbits, persist for billions of years
without merging. The stabilizing mechanism in every such case is
orbital motion: angular momentum.

\textbf{[Tension identified, not resolved.]} Real angular momentum
needs mass ($L=r\times mv$) and inertia, motion that persists between
moments rather than being reset by the current force alone. This
paper's own motion is deliberately mass-free --- the reason mass
segregation does not reproduce here, already a settled point above
--- and deliberately overdamped, velocity set directly by force with
no memory of the step before. A literal transplant of angular
momentum would undo both of those deliberate choices at once.

\textbf{[Tested, did not work; problem remains open.]} A proxy
attempt: a tangential force built from each atom's own $\hat n$,
already present and otherwise unused for force direction, added to
the radial pull. The first check looked promising --- the largest
group shrank from $97.5\%$ to $7.5\%$ of the population as the
tangential term strengthened. A second, independent check showed why
that was misleading: the system was not settling at all at the
stronger values, only still moving (last-step size $\approx0.40$--$0.44$,
against $\approx0.0003$ for the undisturbed case). An overdamped
system has nothing to damp a tangential force the way real inertia
damps orbital perturbation, so the raw proxy trades collapse for
motion that never settles, rather than producing anything stable.
The fix needs a genuine damping mechanism, not yet designed. This is
the most current, least resolved thread in this record.


\addcontentsline{toc}{section}{Aside: A Note Parked Here for the Companion Paper}

This belongs beside \emph{Pattern Arithmetic}'s own Thesis, not
here --- that paper is frozen for now, so it is held in this one
until it migrates.

Not every axiom is earned the same way. Some are forced: $\pi$
($3.14159265\ldots$) and $e$ ($2.71828182\ldots$) carry no freedom at
all, falling out of their own definitions identically regardless of
who does the deriving --- discovered, not invented, in the strongest
sense available. Some are genuinely invented, the way physics' free
parameters are: the fine-structure constant, the particle masses,
Planck's constant's numerical value --- structure discovered,
specific value measured, with no known reason it could not have come
out otherwise. And some, like the companion paper's own $r\le1$, are
neither: a normalization, a choice of scale, cheap to invent because
no structural weight rides on the specific number --- rescale it and
nothing about the theory's shape changes.

The test that actually sorts these, worth stating plainly: not
whether a rule was discovered or invented, but whether a different
choice would have produced a different theory. $\pi$ fails this test
by having no alternative to compare against. $r\le1$ fails it by
mattering only as a scale. Mix as phasor summation, or Gather's
idempotence, pass it --- a different choice there is a different
algebra, not the same one in different clothes. This is the
reasoning underneath the companion paper's own
line, that it ``states its own rules as axioms rather than
discovering them''; it is the same standard this paper's own closure
axioms, wherever they are eventually stated, should be held to.

\end{document}
